\documentclass[11pt,a4paper,sans]{moderncv}

\usepackage{xcolor}

% Theme choice
\moderncvstyle{classic} % options: 'classic', 'casual', 'oldstyle', 'fancy', 'banking'
\moderncvcolor{green}    % options: blue, orange, green, red, purple, grey, black

% Adjust margins
\usepackage[scale=0.9]{geometry}


% Width of the left (years) column — do this once for all sections
\setlength{\hintscolumnwidth}{3.6cm}
\recomputelengths

% Make the left column text dark & bold (and upright, not italic)
% \renewcommand*{\hintfont}{\bfseries\color{black}}


% Personal data
\renewcommand*{\namefont}{\fontsize{50}{60}\bfseries\rmfamily}
\name{Rajesh}{Kumar}
\title{Assistant Professor of Physics\\\textsc{Sido Kanhu Murmu University, Dumka}}
\email{kr.rajesh.phy@gmail.com}
\phone[mobile]{+91~7352949118}
\social[github]{rajeshphy.github.io}
\address{Dumka, Jharkhand}{India}
\extrainfo{Born: 11 April 1995}
\photo[64pt][0.3pt]{profile.jpeg}

\begin{document}
\makecvtitle

\section{Research Interests}
\cvitem{}{Quantum Mechanics, Spectral Methods, Computational Physics}

\section{Education}
\cventry{2016--2018}{M.Sc. Physics}{IIT Kanpur}{}{}{CGPA: 7.4/10}
\cventry{2013--2016}{B.Sc. Physics (Hons.)}{St. Xavier’s College, Ranchi}{}{}{80\%}
\cventry{2011--2013}{Intermediate}{Jharkhand Academic Council}{}{}{71\%}
\cventry{2011}{Matriculation}{ICSE Board}{}{}{88\%}

\section{Research Publications}

\cvitem{2024}{R. Kumar, R.K. Yadav, A. Khare, \textbf{Rationally extended harmonic oscillator potential, isospectral family and the uncertainty relations}, Annals of Physics, 463 (2024): 169623. 
\href{https://doi.org/10.1016/j.aop.2024.169623}{\textcolor{blue}{DOI}}}

\cvitem{2025}{R. Kumar, R.K. Yadav, A. Khare, \textbf{Rational extension of anisotropic harmonic oscillator potentials in higher dimensions}, Annals of Physics, 479 (2025): 170045. 
\href{https://doi.org/10.1016/j.aop.2025.170045}{\textcolor{blue}{DOI}}}

\cvitem{2025}{R.K. Yadav, R. Kumar, A. Khare, \textbf{Rational Extension of Quantum Anisotropic Oscillator Potentials with Linear/Quadratic Perturbations}, arXiv:2504.08236 (2025). 
\href{https://arxiv.org/abs/2504.08236}{\textcolor{blue}{arXiv}}}



\section{Research \& Professional Experience}
\cventry{Oct 2022--Present}{Assistant Professor}{Model College Dumka \& SKMU University}{}{}{Teaching PG Physics; mentoring computational and theoretical projects.}
\cventry{Nov 2020--Oct 2022}{Assistant Professor}{Deoghar College, SKMU University}{}{}{Taught undergraduate and postgraduate physics courses.}
% \cventry{Apr 2020--Nov 2020}{Subject Matter Expert}{Six Red Marbles Learning Pvt. Ltd., Chennai}{}{}{Authored physics/mathematics solutions for international textbooks (LaTeX, JSX).}
\cventry{Mar 2019--Mar 2020}{Project Scientist--B}{INCOIS, Hyderabad}{}{}{Analyzed ocean wave spectra with Python \& shell scripts; mesh generation with Gmsh.}
\cventry{Sep 2018--Feb 2019}{Junior Research Fellow}{Nanophotonics Lab, IIT-ISM Dhanbad}{}{}{Modeled surface plasmon resonance via ellipsometry (theoretical \& computational).}

\section{Short Courses}
\cventry{Sep 2024}{Planetary Science (Online)}{CSSTEAP (UN-affiliated), PRL Ahmedabad}{}{}{Completed short course on planetary science conducted by PRL under CSSTEAP.}

\section{Achievements}
\cventry{2018}{GATE Physics}{}{}{}{All India Rank 484}
\cventry{2017}{CSIR-JRF (Physics)}{}{}{}{All India Rank 218}
\cventry{2016}{National Graduate Physics Examination (NGPE)}{}{}{}{State Topper, Bihar \& Jharkhand (Conducted by IAPT)}
\cventry{2011, 2013}{Block Topper}{}{}{}{Class 10 \& 12, Silli, Ranchi}


\section{Technical Skills}
\cvitem{Programming:}{Python, C, C++, Java, Shell scripting, LaTeX.}
\cvitem{Web Development:}{Django, HTML, CSS, JS, Jekyll.}
\cvitem{Tools:}{Gmsh, Tkinter GUI, Git.}


\end{document}